% Standalone LaTeX theory block for Spectral-Spike regularization.
% This version deliberately avoids the TD-operator recursion and proves
% Bellman-loss decrease directly under a frozen target.

\section{Theory: Direct Bellman-Loss Improvement from a Spectral Spike}
\label{sec:direct-spike-theory}

For a minibatch $\{(h_i,y_i)\}_{i=1}^N$, let
\begin{equation}
    \mathcal B(W)
    :=\frac{1}{2N}\sum_{i=1}^N\delta_i(W)^2,
    \qquad
    \delta_i(W):=Q_W(h_i)-y_i,
    \label{eq:fixed-target-bellman-loss}
\end{equation}
where the target $y_i$ is stop-gradient during the analyzed critic update.  For
clarity, consider a two-layer critic
\begin{equation}
    Q_W(h)=\sum_{k=1}^m o_k\varphi(w_k^\top h),
    \qquad W=[w_1,\ldots,w_m]^\top,
    \label{eq:critic}
\end{equation}
with fixed output weights $o_k$ during the directional argument.  Let
$(s_1,\hat u,\hat v)$ be a top singular triplet of $W$ and define
\begin{equation}
    D:=\hat u\hat v^\top,
    \qquad \|D\|_F=1.
    \label{eq:spike-direction}
\end{equation}

\paragraph{Assumptions.}
We use four standard local assumptions.
\begin{enumerate}
    \item[(A1)] \textbf{Frozen target.}  The values $y_i$ in
    \eqref{eq:fixed-target-bellman-loss} do not depend on $W$ during the
    analyzed update.
    \item[(A2)] \textbf{Directional smoothness.}  There is $L_B<\infty$ such
    that, for all sufficiently small $\alpha$,
    \begin{equation}
        \mathcal B(W+\alpha D)
        \le \mathcal B(W)
        +\alpha\langle\nabla_W\mathcal B(W),D\rangle
        +\frac{L_B}{2}\alpha^2.
        \label{eq:directional-smoothness}
    \end{equation}
    \item[(A3)] \textbf{Supercritical spectral alignment.}  The top right
    singular vector obeys the signal-plus-noise decomposition
    \begin{equation}
        \hat v=\tau v^\star+\sqrt{1-\tau^2}\,v_\perp,
        \qquad \tau\ge\tau_0>0,
        \label{eq:bbp-alignment}
    \end{equation}
    with $v_\perp\perp v^\star$ and unit vectors $v^\star,v_\perp$.  This is
    the usual consequence of operating above the BBP transition.
    \item[(A4)] \textbf{Useful signal and centered orthogonal fluctuation.}
    With $g_i$ and $\xi_i$ defined in Lemma~\ref{lem:spike-feature},
    \begin{equation}
        -\frac1N\sum_{i=1}^N\delta_i g_i\ge c_g>0.
        \label{eq:useful-signal}
    \end{equation}
    Conditional on the current critic, residuals, and signal components,
    $X_i:=\delta_i\xi_i$ are independent, centered, and
    $\sigma^2$-sub-Gaussian.  Conditional independence may be replaced by any
    martingale or mixing assumption yielding the same concentration bound.
\end{enumerate}

\begin{lemma}[Spike-induced tangent feature]
\label{lem:spike-feature}
Let
\begin{equation}
    \psi_i:=\left.\frac{\mathrm d}{\mathrm d\alpha}
    Q_{W+\alpha D}(h_i)\right|_{\alpha=0}.
\end{equation}
Write $h_i=T_i v^\star+h_i^\perp$, where
$(v^\star)^\top h_i^\perp=0$.  Then
\begin{equation}
    \psi_i=S_i(W)\hat v^\top h_i,
    \qquad
    S_i(W):=\sum_{k=1}^m o_k\hat u_k\varphi'(w_k^\top h_i),
    \label{eq:exact-tangent-feature}
\end{equation}
and, under \eqref{eq:bbp-alignment},
\begin{equation}
    \psi_i=\tau g_i+\sqrt{1-\tau^2}\,\xi_i,
    \label{eq:signal-noise-feature}
\end{equation}
where
\begin{equation}
    g_i:=S_i(W)T_i,
    \qquad
    \xi_i:=S_i(W)v_\perp^\top h_i^\perp.
\end{equation}
\end{lemma}

\begin{theorem}[Direct Bellman-loss decrease above the BBP threshold]
\label{thm:direct-bellman-decrease}
Suppose (A1)--(A4) hold.  For any $\delta\in(0,1)$, define
\begin{equation}
    \varepsilon_N(\delta)
    :=\sigma\sqrt{\frac{2\log(2/\delta)}{N}},
    \qquad
    \mu_N(\delta)
    :=\tau_0c_g-\sqrt{1-\tau_0^2}\,\varepsilon_N(\delta).
    \label{eq:descent-margin}
\end{equation}
If $\mu_N(\delta)>0$, then, with probability at least $1-\delta$,
every step size
\begin{equation}
    0<\alpha\le\frac{\mu_N(\delta)}{L_B}
    \label{eq:valid-step}
\end{equation}
satisfies
\begin{equation}
    \boxed{
    \mathcal B(W+\alpha\hat u\hat v^\top)
    \le
    \mathcal B(W)-\frac{\alpha\mu_N(\delta)}{2}.}
    \label{eq:main-decrease}
\end{equation}
Thus a supercritical spectral spike supplies a certified descent direction for
the fixed-target Bellman loss whenever its task-aligned signal dominates the
centered orthogonal fluctuation.
\end{theorem}

\begin{corollary}[Constant-order supercritical gain]
\label{cor:constant-gain}
Under the assumptions of Theorem~\ref{thm:direct-bellman-decrease}, if
$N\ge 8\sigma^2(1-\tau_0^2)\log(2/\delta)/(\tau_0^2c_g^2)$, then
\begin{equation}
    \mu_N(\delta)\ge\frac{\tau_0c_g}{2}.
\end{equation}
Consequently, for
$0<\alpha\le\tau_0c_g/(2L_B)$,
\begin{equation}
    \mathcal B(W+\alpha D)
    \le \mathcal B(W)-\frac{\alpha\tau_0c_g}{4}
\end{equation}
with probability at least $1-\delta$.  The guaranteed decrease is therefore
constant order above the BBP threshold rather than vanishing with dimension.
\end{corollary}

\begin{proposition}[The hinge spike loss moves in the analyzed direction]
\label{prop:hinge-direction}
Let
\begin{equation}
    \mathcal L_{\rm spike}(W):=(T-s_1(W))_+^2,
\end{equation}
where $s_1(W)$ is simple and $s_1(W)<T$.  Then
\begin{equation}
    \nabla_W\mathcal L_{\rm spike}(W)
    =-2(T-s_1(W))\hat u\hat v^\top.
\end{equation}
Hence an isolated gradient step on this regularizer is
\begin{equation}
    W^+=W+2\eta_{\rm sp}(T-s_1(W))\hat u\hat v^\top.
    \label{eq:regularizer-step}
\end{equation}
If the effective step
$\alpha=2\eta_{\rm sp}(T-s_1(W))$ satisfies
\eqref{eq:valid-step}, Theorem~\ref{thm:direct-bellman-decrease} applies to the
regularizer-induced component without changing the SST algorithm.
\end{proposition}

\begin{corollary}[Finite-width robustness]
\label{cor:finite-width}
Suppose the realized directional feature is
$\psi^{(m)}=\psi+r_m$ and
\begin{equation}
    \left|\frac1N\sum_{i=1}^N\delta_i r_{m,i}\right|
    \le\frac{\mu_N(\delta)}{2}.
    \label{eq:finite-width-condition}
\end{equation}
Then, with the same probability as in
Theorem~\ref{thm:direct-bellman-decrease},
\begin{equation}
    \mathcal B(W+\alpha D)
    \le \mathcal B(W)-\frac{\alpha\mu_N(\delta)}{4}
\end{equation}
for every $0<\alpha\le\mu_N(\delta)/(2L_B)$.
\end{corollary}

\appendix
\section{Proofs for the Direct Bellman-Loss Theory}
\label{app:direct-proofs}

\subsection{Proof of Lemma~\ref{lem:spike-feature}}

The $k$th row of $W+\alpha D$ is
$w_k+\alpha\hat u_k\hat v$.  Therefore
\begin{align}
    Q_{W+\alpha D}(h_i)
    &=\sum_{k=1}^m o_k
      \varphi\!\left((w_k+\alpha\hat u_k\hat v)^\top h_i\right),\\
    \left.\frac{\mathrm d}{\mathrm d\alpha}
      Q_{W+\alpha D}(h_i)\right|_{\alpha=0}
    &=\sum_{k=1}^m o_k\hat u_k
      \varphi'(w_k^\top h_i)\hat v^\top h_i\\
    &=S_i(W)\hat v^\top h_i.
\end{align}
This proves \eqref{eq:exact-tangent-feature}.  Next substitute
$\hat v=\tau v^\star+\sqrt{1-\tau^2}v_\perp$ and
$h_i=T_i v^\star+h_i^\perp$.  Orthogonality gives
\begin{equation}
    \hat v^\top h_i
    =\tau T_i+\sqrt{1-\tau^2}\,v_\perp^\top h_i^\perp.
\end{equation}
Multiplying by $S_i(W)$ yields
\begin{equation}
    \psi_i
    =\tau S_i(W)T_i
     +\sqrt{1-\tau^2}S_i(W)v_\perp^\top h_i^\perp
    =\tau g_i+\sqrt{1-\tau^2}\,\xi_i,
\end{equation}
which proves \eqref{eq:signal-noise-feature}.

\subsection{Concentration of the orthogonal fluctuation}

\begin{lemma}[Conditional concentration]
\label{lem:conditional-concentration}
Under (A4), with probability at least $1-\delta$,
\begin{equation}
    \left|\frac1N\sum_{i=1}^N\delta_i\xi_i\right|
    \le\varepsilon_N(\delta).
    \label{eq:concentration-event}
\end{equation}
\end{lemma}

\begin{proof}
Condition on the current critic, residuals, and signal components.  By (A4),
$X_i=\delta_i\xi_i$ are independent, centered, and
$\sigma^2$-sub-Gaussian.  Hence
\begin{equation}
    \Pr\!\left(
      \left|\frac1N\sum_{i=1}^N X_i\right|\ge t
      \ \middle|\ W,\{\delta_i,g_i\}_{i=1}^N
    \right)
    \le 2\exp\!\left(-\frac{Nt^2}{2\sigma^2}\right).
\end{equation}
Setting $t=\sigma\sqrt{2\log(2/\delta)/N}$ makes the right-hand side
equal to $\delta$.  Removing the conditioning proves
\eqref{eq:concentration-event}.
\end{proof}

\subsection{The spike direction has negative Bellman directional derivative}

\begin{lemma}[Descent certificate]
\label{lem:descent-certificate}
Under the event \eqref{eq:concentration-event},
\begin{equation}
    \langle\nabla_W\mathcal B(W),D\rangle
    \le-\mu_N(\delta).
    \label{eq:gradient-certificate}
\end{equation}
\end{lemma}

\begin{proof}
By (A1), the targets are constant with respect to $W$.  The chain rule and
Lemma~\ref{lem:spike-feature} give
\begin{align}
    \langle\nabla_W\mathcal B(W),D\rangle
    &=\left.\frac{\mathrm d}{\mathrm d\alpha}
       \mathcal B(W+\alpha D)\right|_{\alpha=0}\\
    &=\frac1N\sum_{i=1}^N\delta_i\psi_i\\
    &=\tau\frac1N\sum_{i=1}^N\delta_i g_i
      +\sqrt{1-\tau^2}\frac1N\sum_{i=1}^N\delta_i\xi_i.
\end{align}
By \eqref{eq:useful-signal}, the first term is at most $-\tau c_g$.
On the concentration event, the second term is at most
$\sqrt{1-\tau^2}\varepsilon_N(\delta)$.  Therefore
\begin{equation}
    \langle\nabla_W\mathcal B(W),D\rangle
    \le-\tau c_g+\sqrt{1-\tau^2}\varepsilon_N(\delta).
\end{equation}
For $\tau\ge\tau_0$, the function
$f(\tau)=\tau c_g-\sqrt{1-\tau^2}\varepsilon_N(\delta)$ is increasing on
$[0,1]$.  Consequently
\begin{equation}
    \tau c_g-\sqrt{1-\tau^2}\varepsilon_N(\delta)
    \ge\tau_0c_g-\sqrt{1-\tau_0^2}\varepsilon_N(\delta)
    =\mu_N(\delta),
\end{equation}
which proves \eqref{eq:gradient-certificate}.
\end{proof}

\subsection{Proof of Theorem~\ref{thm:direct-bellman-decrease}}

By Lemma~\ref{lem:conditional-concentration}, the event used in
Lemma~\ref{lem:descent-certificate} holds with probability at least
$1-\delta$.  On this event, (A2) and \eqref{eq:gradient-certificate} imply
\begin{equation}
    \mathcal B(W+\alpha D)
    \le\mathcal B(W)-\alpha\mu_N(\delta)+\frac{L_B}{2}\alpha^2.
\end{equation}
If $0<\alpha\le\mu_N(\delta)/L_B$, then
$(L_B/2)\alpha^2\le\alpha\mu_N(\delta)/2$.  Thus
\begin{equation}
    \mathcal B(W+\alpha D)
    \le\mathcal B(W)-\frac{\alpha\mu_N(\delta)}{2},
\end{equation}
as claimed.

\subsection{Proof of Corollary~\ref{cor:constant-gain}}

The sample-size condition is equivalent to
\begin{equation}
    \sqrt{1-\tau_0^2}\,\sigma
    \sqrt{\frac{2\log(2/\delta)}{N}}
    \le\frac{\tau_0c_g}{2}.
\end{equation}
Substitution into \eqref{eq:descent-margin} yields
$\mu_N(\delta)\ge\tau_0c_g/2$.  The permitted step size and loss decrease
then follow directly from Theorem~\ref{thm:direct-bellman-decrease}.

\subsection{Proof of Proposition~\ref{prop:hinge-direction}}

When $s_1(W)$ is simple, the standard singular-value differential is
\begin{equation}
    \nabla_Ws_1(W)=\hat u\hat v^\top.
\end{equation}
For $s_1(W)<T$, differentiation of
$(T-s_1(W))^2$ gives
\begin{equation}
    \nabla_W\mathcal L_{\rm spike}(W)
    =-2(T-s_1(W))\nabla_Ws_1(W)
    =-2(T-s_1(W))\hat u\hat v^\top.
\end{equation}
Subtracting $\eta_{\rm sp}$ times this gradient yields
\eqref{eq:regularizer-step}.  The last statement follows by identifying its
coefficient with $\alpha$ in Theorem~\ref{thm:direct-bellman-decrease}.

\subsection{Proof of Corollary~\ref{cor:finite-width}}

The realized directional derivative is
\begin{align}
    \frac1N\sum_{i=1}^N\delta_i\psi_i^{(m)}
    &=\frac1N\sum_{i=1}^N\delta_i\psi_i
      +\frac1N\sum_{i=1}^N\delta_i r_{m,i}\\
    &\le-\mu_N(\delta)+\frac{\mu_N(\delta)}{2}
     =-\frac{\mu_N(\delta)}{2}.
\end{align}
Directional smoothness therefore gives
\begin{equation}
    \mathcal B(W+\alpha D)
    \le\mathcal B(W)-\frac{\alpha\mu_N(\delta)}{2}
      +\frac{L_B}{2}\alpha^2.
\end{equation}
For $0<\alpha\le\mu_N(\delta)/(2L_B)$, the quadratic term is at most
$\alpha\mu_N(\delta)/4$, proving
\begin{equation}
    \mathcal B(W+\alpha D)
    \le\mathcal B(W)-\frac{\alpha\mu_N(\delta)}{4}.
\end{equation}

\subsection{Scope of the result}

The loss in \eqref{eq:fixed-target-bellman-loss} is the empirical squared
Bellman/TD loss for a frozen target, not a global policy-evaluation theorem.
The result neither invokes the TD recursion nor claims that the moving-target
loss decreases monotonically across all training iterations.  Its precise claim
is local and directly verifiable: above the BBP threshold, the top singular
direction is a Bellman-loss descent direction with a non-vanishing margin when
the aligned signal is useful and the orthogonal component is centered.
